\documentclass[10pt,a4paper]{amsart}
\usepackage[margin=19mm]{geometry}
\usepackage{amsmath,amssymb,mathtools}
\usepackage[scr=rsfs,cal=euler]{mathalfa}
\usepackage{xfrac}
\usepackage[unicode,hidelinks,bookmarks=false]{hyperref}
\newcommand{\ie}{i.e.\ }
\newcommand{\cf}{cf.\ }
\newcommand{\resp}{respectively\ }
\newcommand{\Subsection}{\subsection}
\providecommand{\nobreakdash}{\nobreak\hbox{-}}
\setlength{\emergencystretch}{2em}
\multlinegap0pt
\usepackage{amssymb}
\usepackage[shortlabels]{enumitem}
\usepackage{float}
\usepackage{xcolor}
\usepackage{csquotes}
\usepackage{tikz}
\usepackage{caption}
\newtheorem{lemma}{Lemma}
\title{On Spectral Properties of Lanzhou Matrix of Graphs}
\author{Madhumitha K V, Harshitha A, Swati Nayak, Sabitha D'Souza}
\date{Source: arXiv:2512.19404v1 (CC BY 4.0)}
\begin{document}
\maketitle

\noindent\emph{Source and licence.} On Spectral Properties of Lanzhou Matrix of Graphs. Version arXiv:2512.19404v1 (CC BY 4.0), distributed by arXiv under the Creative Commons Attribution 4.0 International licence, http://creativecommons.org/licenses/by/4.0/, as recorded in the arXiv licence metadata for this exact version and shown on the abstract page https://arxiv.org/abs/2512.19404v1. Full licence notice: \texttt{licenses/arxiv-2512.19404-CC-BY-4.0.txt}.\par\smallskip

\noindent\textit{Editorial scope.} Counted unit: the article's lemma ddd (source0.tex lines 944--949). The article's complete original proof of it is delivered with it (source0.tex lines 950--1003, one \texttt{proof} environment of 54 lines). No mathematics is written, repaired or re-derived: the statement and the proof are the article's own lines, in the article's own notation, and no retained line is substituted or re-flowed.  No further source-local context is needed: every cross-reference of every retained line resolves inside the two delivered spans.  Nothing was omitted from any retained span, so the package carries no omission record.  No retained line contains a citation, so the wrapper carries no bibliography and no bibliography entry.  No retained line contains a figure environment, an image-inclusion command or drawing code, so no image is delivered and the package declares no asset dependency.  Editorial formatting only, in addition: an AMS \texttt{amsart} preamble that restates the article's own declarations and macros used by the retained lines, namely the environments \texttt{lemma} on separate counters, together with the standard packages those lines need.  The retained environments print in this document as Lemma 1 in order of appearance, whereas the article prints the same items with the numbering of its own full document.  The complete native source of the article as distributed in the arXiv e-print for this exact version is bundled unchanged as \texttt{sources/191471/source0.tex}.
\begin{lemma}\label{ddd}
Let $\Gamma$ be a simple, undirected graph. Let $A_w(\Gamma)$ be the real, non-negative weighted adjacency matrix associated with the graph $\Gamma.$ Then the eigenvalues of $A_w(\Gamma)$ are symmetric about the origin (if $\lambda$ is an eigenvalue of $A_w(\Gamma)$ with multiplicity $k$ then $-\lambda$ is also an eigenvalue of $A_w(\Gamma)$ with multiplicity $k$) if and only if $A_w(\Gamma)$ is of the form  $$\begin{bmatrix}
0 & B\\
B^T & 0\\
\end{bmatrix}.$$
\end{lemma}

\begin{proof}
Let $A_w(G)$ be the real, non-negative weighted adjacency matrix associated with the graph $G.$ 
	
Let $A_w(G)=\begin{bmatrix}
0 & B\\
B^T & 0\\
\end{bmatrix}.$ To prove eigenvalues of $A_w(G)$ are symmetric about the origin,	let $\mathbf{x}=[\mathbf{y}\ \mathbf{z}]^{T}$ be an eigenvector of $A_{w}(G)$ with $\gamma$ as the eigenvalue. Then $$\begin{bmatrix}
		0 & B\\
		B^{T} & 0\\
	\end{bmatrix}\begin{bmatrix}
		\mathbf{y}\\\mathbf{z}
	\end{bmatrix}=\gamma \begin{bmatrix}
		\mathbf{y}\\\mathbf{z}
	\end{bmatrix}.$$
	This implies, $B\mathbf{z}=\gamma \mathbf{y}$ and $B^T \mathbf{y}=\gamma \mathbf{z}.$
	
	Define $\mathbf{x'}=\begin{bmatrix}
		\mathbf{y} \\\mathbf{-z}
	\end{bmatrix}\neq 0.$
	
	$$\begin{bmatrix}
		0 & B\\
		B^{T} & 0\\
	\end{bmatrix}\begin{bmatrix}
		\mathbf{y}\\\mathbf{-z}
	\end{bmatrix}=\begin{bmatrix}
		-B\mathbf{z}\\B^{T}\mathbf{y}
	\end{bmatrix}=\begin{bmatrix}
		-\gamma \mathbf{y}\\\gamma
        \mathbf{z}
	\end{bmatrix}=-\gamma \begin{bmatrix}
		\mathbf{y}\\\mathbf{-z}
	\end{bmatrix}.$$
	
	This implies $-\gamma$ is also an eigenvalue with eigenvector $[\mathbf{y}\ \mathbf{-z}]^{T}.$

    Conversely, let the eigenvalues of $A_w(G)$ be symmetric about the origin.
	
	Let \([A_{w}(G )^{k}]_{ij}\) denote the \((i,j)\)-th entry of the matrix \(A_{w}(G )\) raised to the \(k\)-th power.
	This implies $-\gamma $ is also an eigenvalue with eigenvector $[y\ -z]^{T}.$
	
	
	
	We have
	$$[A_{w}(G )^{k}]_{ij}=\sum\limits_{\substack{closed \ walks \\ of \ length\  k\\  starting\\  at\  i\  and \\ ending\ at\  j}}\prod\limits_{\substack{edges\ in \\ the \ walk\\ i\leq u<v\leq j}}[A_{w}(G )]_{uv}.$$
Since the eigenvalues of $A_w(G)$ are symmetric about the origin, we have
$trace[A_{w}(G)^k]=0,$ for all odd $k.$ 
Since $A_w(G)$ is non-negative, there is no cycle of odd length in $G.$ 

$\implies$ $G$ is bipartite. With a proper labeling of $V(G)$, we can find a a non-trivial partition of $V(G)$ such that $$A_w(G)=\begin{bmatrix}
0 & B\\
B^T & 0\\
\end{bmatrix}.$$ 
\end{proof}

\end{document}
